\Needspace{18\baselineskip}
\section{Scope, notation, and the status of the formalization}

The abstract part of the development follows the proof in Appendix~F of
\cite{paper}. Its completed mathematical core is the soundness of counting
simulation, corresponding to Proposition~37. The remaining argument is expressed
through a structure of abstract specimens. Given that structure and its laws,
Lean proves both directions of the simulation construction and their contextual
consequences. Constructing this structure from the paper's actual hypernets is
still a separate task. The concrete graph part now supplies finite interfaced
graphs, certified shallow plugging, renaming and path invariance, a typed boxed
hierarchy, and addressed deep replacement. These constructions and their
regression examples are explained after the abstract development. The context
calculus now also proves exact hole inventories, necessary-and-sufficient
freshness, and exact consumption of selected hole labels at arbitrary box depth.
Recursive isomorphism now identifies independently renamed box contents and
respects paths, hole observations, and deep replacement. Sites now transport across
isomorphic hosts, and substitution respects simultaneous host and replacement
renaming, including the exact freshness condition. Nested substitutions now
associate at arbitrary box depth, with corolla unit laws and a sharp final
context-validity criterion even when an intermediate collision is later removed.

The concrete development now also defines globally focus-free hypernets,
focussed hypernets, and the hole-free state space. The boundary argument of
Lemma~39(1) is proved for finite sequences of addressed focus-free fillings;
see Section~\ref{sec:concrete-focus}. The five search rules are now defined on underlying graph coordinates, with
determinism, boundary laws, and preservation under shallow filling proved
(Section~\ref{sec:search}). Explicit focus insertion now realizes typed
coordinates as actual states, preserving holes and boundary observations, and
contracting each inserted focus recovers original incidences
(Section~\ref{sec:realization}). General erasure and the focus-preserving erasure-then-insertion isomorphism are
now proved for arbitrary closed focussed states. Erasure-based search on actual
states is deterministic up to focussed isomorphism and has the expected search
boundary laws (Section~\ref{sec:state-search}). Both focus round trips and
source/target renaming invariance are proved; search and boundary observations
descend to literal-label state isomorphism classes, with equality-based search
determinism. Search also respects recursive renaming inside boxes and descends
to recursive state classes. The map from literal to recursive classes preserves
and lifts search executions with exact lengths, including successful search
termination. Canonical roots and concrete rootedness are now defined; search
preserves rootedness, and concrete rooted focus sources have operation paths
from the unique input. An incoming path containing a non-operation edge rules
out rootedness (Section~\ref{sec:rooted-path-obstruction}). Concrete refocusing
is defined and proved for search; Section~\ref{sec:uam-structure} now proves
the complete UAM structural theorem. One-way hypernets and the single-hole
obstruction of Lemma~26(2) are now proved, using extraction of actual operation
paths through shallow substitution (Section~\ref{sec:one-way-obstruction}).
Simultaneous shallow expansion now has proved graph invariants. Concrete
pre-templates, output-closure, and shallow specimens provide actual filled
states and a multi-hole proof of Lemma~26(3) for this representation
(Section~\ref{sec:shallow-specimens}). Initial/final agreement now applies
directly to these specimen states. Every finite simultaneous shallow filling
now has a finite ordinary addressed-substitution realization, and completed
schedules give isomorphic results (Section~\ref{sec:sequential-expansion}).
Independent deep substitutions now have constructed residual sites and a
recursive interchange isomorphism. A quotient of hosts with selected sites
supports substitution and sharp context validity (Section~\ref{sec:deep-interchange}).
Recursive deep specimens now supply actual filled states, with derived
focus freedom and the Lemma~26(3), boundary, and jointly rooted context-path
conclusions (Section~\ref{sec:deep-specimens}). Recursive fillings now preserve
context search with exact execution lengths. Search at host coordinates reflects
to context search away from shallow hole boundaries; invariant regions support
reflection of entire runs. These coordinate executions also realize paired
executions and rooted endpoints on actual state classes
(Section~\ref{sec:expansion-search}). General erasure now commutes with
expansion, identifying these realizations with the original deep specimen
states. Concrete entering/exiting predicates agree with the coordinate boundary
test, output closure and joint rootedness imply non-exiting, and settled
specimens have synchronized actual search transitions to rooted endpoints at
a remote coordinate (Section~\ref{sec:specimen-traversal}). Section~\ref{sec:concrete-specimen-system} now constructs the named-state interface and safety/robustness bridges.
Section~\ref{sec:tuple-safety} proves compatibility with finite original-hole
tuples and defines concrete quasi-specimens, input-safety, and robustness.
Contextual search adequacy is proved below. Concrete substitution and compute
transitions now assemble into the full UAM, with mode discipline, disjoint
transition families, and final normality proved (Section~\ref{sec:full-uam}).

\Needspace{18\baselineskip}
\subsection{Paper-to-Lean proof dependency map}
The target is the proof of Theorem~6 through Propositions~40, 37, and 38.
The following map separates checked arguments from their missing concrete
interpretations, rather than treating graph infrastructure as progress on every
part of that proof.
\begin{enumerate}
\item \textbf{Proposition~37: checked.}\par
\texttt{Machine.CountingSimulation.sound} proves counting-simulation soundness
for a deterministic machine and a reasonable cost triple. Its proof follows
the quantitative argument, including the eventual-stuck alternative.
\item \textbf{Proposition~40: concrete forward simulation checked.}\par
\texttt{SpecimenSystem.countingSimulation} and its converse construct counting
simulations from the abstract specimen laws, safety, and robustness. The abstract laws have concrete named-state and class-valued instances. Section~\ref{sec:uam-classes} proves literal full-class determinism and the concrete forward counting simulation using class safety and robustness. Section~\ref{sec:class-converse} completes the concrete converse counting simulation and two-sided refinement. Section~\ref{sec:concrete-observations} completes those observations and both contextual clauses for the fixed-depth machine. Operational transfer to the unbounded machine remains open.
\item \textbf{Lemma~39(1), Proof~F.4: boundary argument checked.}\par
\texttt{Focussed.fillings\_boundary\_agreement} proves initial/final agreement
for two finite filling sequences of one focussed hypernet. The proof preserves
the focus mode and input status of its source, as in the paper. Finite original-hole tuples and the recursive representation are now
related in both directions in Section~\ref{sec:tuple-safety}.
Both the shallow and recursive deep specimen representations now have direct
proofs of this boundary conclusion for their actual filled states.
\item \textbf{Lemma~39(2): checked for recursive specimens; point~(3): full-UAM locality checked.}\par
Concrete rootedness and output closure yield the non-operation-path obstruction
of Lemma~26(3) for actual recursive specimen states. An exiting focus would give
a singleton hole path to the focus source, so joint rootedness rules it out,
proving the concrete conclusion of point~(2).

General erasure commutes with simultaneous expansion, preserving the joined
coordinate. Consequently expansion of an already focussed context is
focus-preservingly isomorphic to filling its erasure and then inserting a focus.
Entering/exiting adjacency in the original context is equivalent to the checked
coordinate boundary test. Under the settled hypotheses of point~(3), search
enabled on either specimen side yields synchronized actual-state steps to rooted
endpoints at a remote coordinate. The proof allows recursive endpoint renaming
and uses output closure to exclude exiting backtracking focuses.

Section~\ref{sec:concrete-specimen-system} now packages that coordinate as
an allowed unique-hole specimen with recursive fillings and closed endpoints.
Its full-UAM context-step theorem constructs the abstract locality field.
Original-hole tuple compatibility is proved in Section~\ref{sec:tuple-safety};
contextual search adequacy is established in Section~\ref{sec:contextual-search}.
\item \textbf{Paper \S6.4: concrete transition definition checked.}\par
\texttt{State.Substitution} applies the explicit Fig.~9(f) fragments in a
shallow simple context. \texttt{State.Compute} specifies the required active
operation and focus modes; \texttt{Behaviour} restricts the extrinsic parameter
to this shape. \texttt{UAM.machine} assembles all three families, with proved
final normality and transition-family exclusion. This removes the transition
definition obligation.
\item \textbf{Lemma~24, Lemmas~27--28, and Lemma~3: structural theorem checked.}\par
Section~\ref{sec:uam-structure} proves chronological search revisiting,
stationary prefix replay, substitution rootedness and uniqueness, and
full-UAM determinism and refocusing exactly under the corresponding extrinsic
behaviour conditions. Full executions preserve rootedness.
\item \textbf{Paper \S\S8.1--8.3: tuple compatibility and definitions checked.}\par
\texttt{DeepSpecimen.substitutionTuple} constructs distinct, complete finite
original-hole tuples with related closed entries. \texttt{TupleFilling.specimen}
reconstructs a recursive specimen with the same filled graphs. Concrete
quasi-specimens retain unconditional witness existence and mode equality;
input-safety is unrestricted by rootedness, and robustness retains its
unrooted alternative. The rooted quasi and robustness consequences are proved.
Section~\ref{sec:concrete-specimen-system} proves the safety and robustness bridges for the constructed named-state \texttt{SpecimenSystem}.
\item \textbf{Proposition~38: logical lifting checked; initial bridge open.}\par
\texttt{Observations.contextual\_of\_lifting} proves the logical implication
once initial pairs belong to the lifting. Concrete initial focus attachment is
now constructed, including initial rooted specimens obtained by relocation.
The observation interface and its arbitrary-fragment witnesses must still
discharge the \texttt{initial\_lift} parameter.
\item \textbf{Theorem~6: conditional, not complete.}\par
The two \texttt{contextual\_sufficiency} theorems combine these abstract pieces.
The paper itself assumes determinism and refocusing of the UAM. These remain
the theorem's hypotheses. Section~\ref{sec:uam-structure} now defines their
concrete meaning and proves the full-UAM structural theorem, reducing both
properties exactly to the corresponding extrinsic behaviour conditions.
\end{enumerate}

The typed finite-depth hierarchy and explicit renaming bijections are
representation choices. Outer set-presentation normalization, recursive annotated round trips, and
compatibility across depth bounds are now checked in
Section~\ref{sec:set-presentations}. Section~\ref{sec:paper-encoding} resolves the remaining box-label correspondence
by fresh naming of actual contents, with inverse class maps and exact depth
coherence. The checked counterexample to naive erasure is repaired without
restricting the typed graph class. The existing plugging,
path, inventory, and isomorphism lemmas support that representation; they do
not replace the operational proofs in Appendix~F. The contextual-search dependency is now discharged:
Section~\ref{sec:contextual-search} identifies erasure-based state search with
finite interaction patterns plugged into explicit simple contexts, in both
directions. Opening and closing the selected slot covers arbitrary simple
contexts; output fragments and self-looping attachments are included.
The transitions of \S6.4 are now assembled into a concrete machine in
Section~\ref{sec:full-uam}. Section~\ref{sec:uam-structure} proves substitution
rootedness and uniqueness, stationary replay, and the full-UAM structural laws.
Section~\ref{sec:concrete-specimen-system} closes the successor packaging,
full-UAM locality, and concrete safety/robustness bridges. The remaining
operational connection is transfer to the union of all finite depths. Section~\ref{sec:uam-classes} constructs the deterministic full
class machine and its concrete specimen system, class safety/robustness
bridges, forward counting simulation, and remote-specimen refinement.

Section~\ref{sec:class-converse} supplies the concrete converse quasi and robustness
bridges, reverse counting simulation, two-sided refinement, and equivalence
for constructed initial specimens.

Section~\ref{sec:concrete-observations} constructs template-independent
observations, proves initial lifting, and instantiates both contextual
sufficiency clauses, including arbitrary single fragments and nested holes.
The fixed-depth concrete theorem is complete. The scope audit identifies
operational transfer to the paper's unbounded carrier as the remaining
mathematical obligation; the existing graph encoding does not prove it.
The earlier 94\% planning estimate should not be read as a measured completion
percentage. Its remaining-effort allowance did not explicitly separate this
operational transfer, so a more precise overall percentage is not justified.
The active to-do list is maintained in Appendix~\ref{app:remaining-work}.

There are three distinct kinds of assertion below. A \emph{definition} specifies
a mathematical object or predicate. A \emph{proved statement} is established by
a Lean proof term. An \emph{interface law} is a hypothesis carried by a
structure; the current code uses that hypothesis but does not prove it for
concrete hypernets. The concrete instances proved later in this document discharge those
interface laws for the fixed-depth universal abstract machine.

All code links point to revision \texttt{a3adde9}, rather than to the moving
\texttt{main} branch. Each link highlights the declaration together with its
proof or defining fields. Namespace prefixes in link labels disambiguate
declarations with the same short name. Anonymous Lean \texttt{example}
declarations are explained alongside the named theorems.

\Needspace{17\baselineskip}
\subsection{Conventions}

Fix a set of states $X$ and a machine $M$ on $X$. Write $s\to t$ for its
one-step transition relation, $s\steps{n}t$ for exactly $n$ transitions, and
$\fin(s)$ for successful finality. Natural numbers include zero. A cost
preorder $Q$ is written infix: $n\,Q\,m$. The source count is always the
left argument and the target count the right argument. Write $s\preceq_Q t$
for state refinement, and $s\simeq_Q t$ for mutual state refinement.

The paper's three cost parameters $(Q,Q',Q'')$ are named $(Q,P,S)$ in the
code and in this commentary. Thus $P$ is the auxiliary relation on the
\emph{left} of a simulation candidate, and $S$ is the one on its
\emph{right}. We spell out witnesses rather than relying on a convention for
composition of relations.

No result assumes that all executions terminate. An arbitrary state can be
final, stuck, or capable of further computation. Successful termination is the
only qualitative observation in the present state-refinement definition.
Neither equality of returned values nor equality of final states across two
different source states is required.

\Needspace{17\baselineskip}
\subsection{Dependency order}

The files are organized as follows. Finite-execution facts in
\leanfile{Hypernets/Transition.lean}{Transition.lean} support the quantitative
relations in \leanfile{Hypernets/Quantitative.lean}{Quantitative.lean} and
the soundness argument in
\leanfile{Hypernets/Simulation.lean}{Simulation.lean}. The observation interface
in \leanfile{Hypernets/Contextual.lean}{Contextual.lean} converts state
results into contextual ones. The specimen argument in
\leanfile{Hypernets/LocalReasoning.lean}{LocalReasoning.lean} constructs the
simulation to which soundness applies. Finally,
\leanfile{Hypernets/Examples.lean}{Examples.lean} checks substantive edge
cases, and \leanfile{Hypernets/Audit.lean}{Audit.lean} reports theorem
dependencies. The root file \leanfile{Hypernets.lean}{Hypernets.lean} imports
the development and its examples; it adds no mathematical declaration.

The concrete branch begins with
\leanfile{Hypernets/Graph/Finite.lean}{Graph/Finite.lean} and
\leanfile{Hypernets/Graph/Basic.lean}{Graph/Basic.lean}. Plugging, paths, and
isomorphism support the hierarchy and deep replacement modules. Operation-path
uniqueness imports the counted-execution results, providing a concrete use of the
abstract machine lemmas. Occurrence fibres and provenance then support shallow and deep context closure
in the new context modules. Both branches are imported by the root file.

The exposition follows these dependency orders. Section numbers and theorem
numbers here are local to this document; references to the paper are explicitly
identified as such.
